\subsection{Prospective Follow-up and Exploratory Additional Model}
Four follow-up checks were locked after the original results. Qwen's commit/BM25 stages completed but were not retained when Stage C failed before artifact staging; recovery also failed, so no Qwen follow-up is reported. Mistral and Luna are the completed locked follow-ups. Gemini 2.5 Flash-Lite is exploratory with a matched baseline; no locked verdict depends on it. Stage C is falsified if any locked model crosses a threshold, so its verdict is independent of Qwen.

\emph{Design.} The frozen dense retriever was audited at top-10. A commit-required prompt reused one-poison $k=5$; CAR uses the same 49 exposed Tier~1 trials, item-cluster intervals, and item-level paired inference. BM25 used lower-cased whitespace tokens and corpus-order ties at $k=2,5$. Stage~C used one formal query per item: C0, two clean supports; C1, poison+paired twin; C2, poison+different clean support; C3, poison+distractor; C4, poison alone. The strong composition hypothesis was falsified if any locked model reached $\geq20\%$ adoption in C1/C2 or $\leq50\%$ in C4.

\emph{Results.} The audit reproduced all 640 original retrieval contexts (1920 generation rows across three models) with no poison-rank mismatch (one-poison $k=5$, three-poison $k=5$, bulletin $k=5$, and one-poison $k=2$). Across both tiers at $k=2$, 9 exposed trials had no clean same-item passage and all 9 adopted; 20 exposed trials had at least one clean same-item passage and none adopted. Because the adoption outcomes were known before this audit, this is a consistency check rather than an independent test. In the one-poison $k=5$ Tier~1 contexts, the paired clean twin was present in 47/49 exposed trials; at $k=2$ it was present in 17/26. The $k=2$ exposed set spans 19 items and the nine adopting trials span 9 items. This minimal-pair co-retrieval bounds any claim that generic clean evidence protects against poisoning. Requiring commitment raised adoption from 2/49 to 24/49 for Mistral (49.0\%; item-cluster 95\% CI 33.3--64.4; item-level exact $p=3.05\times 10^{-5}$, $n=31$ items) and from 2/49 to 8/49 for Luna (16.3\%; 7.5--26.5; item-level exact $p=0.0312$, $n=31$ items). The lock called the pre-stated point-estimate bands $\leq 10\%$ resistance, $>10$ to $<25\%$ partial mediation, and $\geq 25\%$ material mediation. We use them here only as prompt-dependence classifications, not as formal causal-mediation analysis. Mistral meets the material band. Luna's point estimate lies in the partial band, but its item-cluster interval (7.5--26.5\%) crosses all three bands. Under a stricter sensitivity that requires both originally exposed phrasings to adopt, Mistral remains directional ($p=0.0312$) while Luna has no positive discordant item ($p=1.0000$). Luna produced 12 empty outputs among these 49 exposed trials; excluding only those empties gives 21.6\% adoption. Pooled clean-control CRR was 83.8\% (Mistral), 86.9\% (Luna), and 84.4\% (Gemini, exploratory). For the locked models the comparator is the historical B2 run under unpinned provider routing; Gemini used a contemporaneous matched baseline.
\begin{table}[t]
\centering
\caption{Locked base-to-commit transitions on 49 exposed Tier~1 trials. N is non-commitment; Other contains transitions not shown separately.}
\label{tab:commit-transitions}
\footnotesize
\begin{tabular}{lrrrrr}
\toprule
Model & A$\to$A & N$\to$A & N$\to$C & N$\to$N & Other \\
\midrule
Mistral & 2 & 22 & 23 & 0 & 2 \\
Luna & 2 & 6 & 26 & 12 & 3 \\
\bottomrule
\end{tabular}
\end{table}

Empty-output audit: Mistral and Qwen had none in B2\_P1/Stage B; Luna had 1/160 in B2\_P1, 4/80 poison-first, none poison-last, and 2/80 in each of C2/C3 (none in C0/C1/C4).
B2 providers were DeepInfra/Parasail/Venice (Mistral), Alibaba (Qwen), and OpenAI (Luna); commit providers were DeepInfra/Parasail/Venice and OpenAI. Closed-book Tier~1 correct counts were 0/70, 2/70, and 0/70; Luna had 41/70 empties.
Stage~C falsified the pre-stated strong composition hypothesis: Mistral adopted in 17/80 C2 trials (21.2\%; exact 95\% CI 12.9--31.8), one trial above the threshold, while Luna (15/80) met every threshold. Both locked models adopted in 0/80 C1 and 80/80 C4 trials. In C1 the paired twin protected integrity by producing non-commitment rather than correct answers: both locked models were 0/80 correct. The no-poison C0 control was 0/80 adoption and 80/80 correct for every model. C2/C3/C4 correct counts were 0/80 for both locked models; Table~\ref{tab:upgrade-controls} reports adoption counts.
\begin{table}[t]
\centering
\caption{Follow-up controls (counts). Commit: attacker adoptions among 49 exposed Tier~1 trials under the commit-required prompt (base prompt: 2, 2, and 3). BM25: adoptions among trials exposed at $k=5$ under the locked corpus-order tie rule. C1--C4: adoptions out of 80 forced-context trials. Gemini 2.5 Flash-Lite is exploratory.}
\label{tab:upgrade-controls}
\footnotesize
\setlength{\tabcolsep}{3.5pt}
\begin{tabular}{lrrrrrr}
\toprule
Model & Commit & BM25 & C1 & C2 & C3 & C4 \\
\midrule
Mistral Small 3.2 & 24/49 & 2/38 & 0 & 17 & 53 & 80 \\
GPT-5.6 Luna & 8/49 & 0/38 & 0 & 15 & 73 & 80 \\
Gemini (expl.) & 21/49 & 12/38 & 1 & 63 & 78 & 80 \\
\bottomrule
\end{tabular}
\end{table}

Under BM25, the poison and its paired clean twin had exactly equal scores in 160/160 queries. Exposure was therefore highly tie-order dependent: under the locked clean-first corpus order, the poison appeared in 2/160 queries at $k=2$ and 38/160 at $k=5$; placing the poison first among equal scores changes those counts to 112/160 and 124/160, respectively. BM25 is therefore a tie-order sensitivity, not evidence of lexical robustness.
\emph{Post-hoc benchmark consequence mapping.} A post-hoc benchmark-consequence mapping decomposed each OSHA MAD into its electrical component and ergonomic component. Of the 18 MAD poisons, 13 consumed only part of the ergonomic component, leaving 0.03--0.45~m beyond the electrical component, and 5 were numerically below the benchmarked electrical component by 0.01--0.17~m. All 7 PRC-024-3 poisons shortened the benchmark no-trip time by 60\%. All 8 Massachusetts DER trip-threshold poisons moved the encoded threshold farther from nominal, widening the benchmark no-trip region. The remaining 2 public exact-value items were return-to-service values and are not assigned a directional consequence class here. All 45 fictional Tier-2 poisons exceeded the record limit by 20\%. 2 of these electrical-component cases (T1-025, T1-027) were among the nine $k=2$ adopted trials. These are benchmark consequence classes, not professional-engineering or physical-safety determinations.
\emph{Post-hoc coverage what-if.} Severity-prioritized D3 gave no consistent benefit: relative to hash coverage, it changed violating trials per model by -1 to +0 of 160 at 40\% and by +1 to +2 at 70\%, with Tier-3 coverage held fixed.

